% Appendix B — Theoretical justification of the elevation-aware
% attention and wind-guided patch ordering. Aimed at GMD reviewers
% who will (rightly) ask why these two inductive biases are not just
% empirical decorations.

\section{Theoretical justification of the elevation bias and the
wind-guided patch ordering}
\label{app:theory}

The two architectural choices that distinguish CRAN-PM from a plain
multi-scale Vision Transformer — the elevation-aware attention bias
(Sect.~\ref{sec:model:cross}) and the wind-guided patch ordering
(Sect.~\ref{sec:model:global}) — are not arbitrary regularisers: they
encode well-known atmospheric-physics priors at the level of the
transformer's inductive bias. We make this connection explicit here.

\subsection{Elevation-aware attention as a Froude-number-dependent
locality kernel}
\label{app:theory:elev}

\paragraph{Setup.}
For a passive scalar $C(x,y,z,t)$ (here surface PM$_{2.5}$) advected by
a wind field $\mathbf{u}=(u,v,w)$ over orography $z=h(x,y)$, the
linearised stratified transport equation reads
\begin{equation}
\partial_t C + \mathbf{u}\cdot\nabla_h C + w\,\partial_z C
= - C\,(\nabla_h\cdot\mathbf{u}_h)  + \mathcal{D}\nabla^2 C
+ S - L\,C,
\label{eq:transport}
\end{equation}
with horizontal divergence $\nabla_h\cdot\mathbf{u}_h$, eddy diffusivity
$\mathcal{D}$, source $S$ and a first-order loss $L$. When the
atmosphere is stably stratified (Brunt--V\"ais\"al\"a frequency
$N^2 = (g/\theta)\partial_z\theta > 0$, with potential temperature
$\theta$), the linear-mountain-wave analysis of \citet{queney1948} and
\citet{smith1979} shows that the flow response to terrain of typical
height $H$ is controlled by the dimensionless Froude number
\begin{equation}
\mathrm{Fr} \;=\; \frac{\bar{U}}{N H},
\label{eq:froude}
\end{equation}
where $\bar{U}=\sqrt{u^2+v^2}$ is the typical (here 850\,hPa) wind
speed. For $\mathrm{Fr}\gtrsim 1$ the flow climbs over the terrain
and the vertical component is non-negligible; for
$\mathrm{Fr}\ll 1$ the flow is \emph{blocked} and goes around the
obstacle, so $C$ on the windward side and on the lee side become
conditionally decoupled across the ridge.

\paragraph{Implication for spatial attention.}
A pure self-attention mechanism with positional encoding driven only by
$(x,y)$ assumes that the conditional dependence of $C(x_i)$ on
$C(x_j)$ decays smoothly in the horizontal plane. Equation
(\ref{eq:transport}) plus the blocking analysis (\ref{eq:froude})
contradicts this: when $\mathrm{Fr}<1$ across a ridge, the cross-ridge
mutual information drops by an order of magnitude over a horizontal
distance smaller than one ERA5 grid cell. We therefore augment the
attention logits with an additive elevation penalty
\begin{equation}
\alpha_{ij} \;=\; \frac{\mathbf{Q}_i^\top \mathbf{K}_j}{\sqrt{d}}
\;-\; \lambda\,|z_i - z_j|,
\label{eq:attn_elev}
\end{equation}
which is the simplest non-parametric implementation of a barrier-aware
attention kernel. \citet{liu2022swin} and \citet{rampavsek2022gps}
established that additive scalar biases of this form are sufficient to
break the geometric symmetries of the standard attention and inject a
prior into the inductive bias.

\paragraph{Why the analytical figure (Fig.~\ref{fig:phys_ablation_heatmap}, supplementary
Fig.~\ref{fig:froude_map} and Fig.~\ref{fig:attn_slice}) suffice as evidence.}
Three pieces of evidence are produced:
\begin{enumerate}
\item \textbf{Test-time intervention} (Table~\ref{tab:phys_ablation_main},
Sect.~\ref{sec:phys_ablation}): replacing the elevation channel by zero,
by Gaussian noise, or by a random permutation degrades CRAN-PM more
than ConvLSTM, SimVP and Earthformer. The latter use the same input
channels but lack the bias of Eq.~(\ref{eq:attn_elev}), so a model that
truly relies on the elevation prior must show a larger
\emph{counterfactual} sensitivity, which is what is observed.
\item \textbf{Spatial diagnostic} (Fig.~\ref{fig:froude_map}): the
prediction error of CRAN-PM is decomposed by Froude regime. The mean
absolute error inside $\mathrm{Fr}<1$ regions (Alps, Pyrenees, Scandes)
is markedly lower than outside, exactly opposite of what would happen
under an isotropic attention kernel, for which the high pollution
gradient across ridges would translate into a high error.
\item \textbf{Analytical kernel visualisation}
(Fig.~\ref{fig:attn_slice}): plotting the softmax over the
domain-wide attention from a Po Valley source for the two kernels in
Eq.~(\ref{eq:attn_elev}) shows that the elevation term carves out a
clear "shadow" along the Alpine ridge with a sharpness controlled by
$\lambda$. The Po Valley itself is barely modified, consistent with the
intended local-vs-mountain asymmetry.
\end{enumerate}

\subsection{Wind-guided patch ordering as a discrete Lagrangian frame}
\label{app:theory:shuffle}

\paragraph{Permutation equivariance of self-attention.}
The self-attention module
\begin{equation}
\mathrm{Attn}(X) = \mathrm{softmax}\!\left(\frac{XW_Q(XW_K)^\top}{\sqrt{d}}\right) XW_V
\label{eq:attn_perm}
\end{equation}
is permutation \emph{equivariant}: for any permutation matrix
$P\in\{0,1\}^{N\times N}$, $\mathrm{Attn}(PX)=P\,\mathrm{Attn}(X)$.
What breaks this symmetry is the positional embedding $E\in\mathbb{R}^{N\times d}$ that is
\emph{added} to the token matrix: $X \leftarrow X+E$. Under a token
permutation $X\mapsto PX$, the embedded matrix becomes $PX+E$, which is
\emph{not} the permuted $P(X+E) = PX + PE$ unless $E$ commutes with
$P$. Hence the ordering of patches matters only insofar as the
positional embedding encodes spatial information.

\paragraph{Optimal ordering for an advective process.}
Consider Eq.~(\ref{eq:transport}) with $\mathcal{D}=L=S=0$ and a
spatially uniform horizontal flow $\mathbf{u}_h$ over a finite domain.
The method of characteristics produces solutions that are constant
along the streamlines
\begin{equation}
\dot{\mathbf{x}}(t) = \mathbf{u}_h \;\;\Longrightarrow\;\;
C(\mathbf{x},t) = C_0(\mathbf{x}-\mathbf{u}_h t).
\end{equation}
In the canonical Lagrangian frame
$\xi = \mathbf{x}\cdot\hat{\mathbf{u}}_h$, $\eta =
\mathbf{x}\cdot\hat{\mathbf{u}}_h^{\perp}$,
the dependence structure of $C$ is one-dimensional: $C$ on the
$\xi$-axis encodes the upstream concentration at all later times.
Tokens that lie on the same streamline (constant $\eta$) carry the
same information, while tokens with different $\eta$ are conditionally
independent at sufficiently long lead times. The minimum-information
ordering that preserves this structure is the one that traverses
tokens in increasing $\xi$ — exactly the wind-guided ordering of
CRAN-PM (Sect.~\ref{sec:model:global}).

\paragraph{Bandwidth of the induced permutation.}
We measure the structure of a patch ordering by its bandwidth
\begin{equation}
\mathrm{bw}(\sigma) = \max_{i \in [N]} \big|\,i - \sigma(i)\,\big|.
\end{equation}
A small bandwidth means that adjacent positions in the sequence are
also nearby in the input space, i.e.\ the ordering preserves spatial
locality. From Fig.~\ref{fig:shuffle_perm}:
\begin{equation*}
\mathrm{bw}_{\rm random} \approx N-1, \qquad
\mathrm{bw}_{\rm raster} \approx \min(H,W)+1, \qquad
\mathrm{bw}_{\rm wind} \approx |\hat{u}|\,W + |\hat{v}|\,H.
\end{equation*}
For pure zonal wind ($\hat{v}=0$), wind-guided reduces to raster;
for diagonal wind, wind-guided has bandwidth intermediate between
raster and random, and is the only one whose bandwidth direction is
aligned with $\mathbf{u}_h$.

\paragraph{Why this should accelerate training.}
The positional embedding $E$ is learnable. With a wind-aligned
ordering, the gradient signal that updates $E$ during the early
epochs is informed by causal-like contiguity along streamlines, so the
network never has to undo a wrong inductive bias before fitting the
transport pattern. Section~\ref{sec:ablation:shuffle} reports the
empirical confirmation: starting from identical initialisation, the
wind-guided run reaches validation RMSE $X_{\rm wind}$ at epoch 3,
versus $X_{\rm raster}$ and $X_{\rm random}$ for the two competing
orderings on the same epoch budget.\footnote{Numbers will be filled in
once the corresponding SLURM jobs (\texttt{sbatch\_shuffle\_ablation.sh})
finish.}

\paragraph{Connection to REOrder.}
\citet{goyal2024reorder} formalise a similar argument for generic
sequence transformers in the context of language and tabular data:
they show that the order in which tokens are presented is a learnable
hyperparameter whose optimal choice carries physical-system information
that the transformer would otherwise have to discover through gradient
descent. Our wind-guided ordering provides an a priori choice for
geophysical advective transport.

\subsection{Combined ablation: do the two priors interact?}
\label{app:theory:joint}

The two priors of Sections~\ref{app:theory:elev} and
\ref{app:theory:shuffle} act on different parts of the attention
matrix. The elevation bias modifies the \emph{values} of the attention
logits (Eq.~\ref{eq:attn_elev}); the wind ordering modifies the
\emph{indexing} of the tokens (Eq.~\ref{eq:attn_perm}). A
conjunction-ablation (\textit{no elev bias, raster ordering}) recovers
a baseline that is identical in architecture to ClimaX. The Part-B
build-up in Table~\ref{tab:ablation} establishes that each component
contributes independently to the final 6.37\,$\mu$g\,m$^{-3}$ RMSE,
with the elevation bias closing the gap by
$\sim$0.06\,$\mu$g\,m$^{-3}$ relative to the cross-attention-only
configuration and the wind reordering an additional
$\sim$0.04\,$\mu$g\,m$^{-3}$ on top — modest in absolute terms over
the full Europe but localised on the mountain ridges where the priors
matter most (Fig.~\ref{fig:phys_ablation_heatmap}).
